% GENERATED FILE. Edit canonical lessons, then run npm run generate:books.
% canonical-source-hash: fnv1a64:6a344be1b2e32ebe
% canonical-lessons: IT-C326-drago, IT-C326-strega, IT-C326-fata, IT-C326-mago, IT-C326-tesoro

\chapter{Draghi e streghe}
\label{ch:it-draghi-e-streghe}
\hlchaptermodality{326}

\begin{chapteropening}
\textbf{By the end of this chapter:} \emph{I can tell a magic tale: the dragon, the witch, the fairy, the wizard and the treasure.}

Complete the last lesson of chapter 326: I can tell a magic tale: the dragon, the witch, the fairy, the wizard and the treasure.
\end{chapteropening}

\section[il drago]{il drago — the dragon}

\label{lesson:IT-C326-drago}



\begin{quote}
Before the new one: say the Italian for the kingdom, then the Italian for to reign.
\end{quote}

\subsection*{You'll want to know: il drago}
\textbf{il drago} — \textquotedblleft{}the dragon\textquotedblright{}.

\begin{quote}
Il drago dorme nella grotta.
\end{quote}

\begin{quote}
\emph{The dragon sleeps in the cave.}
\end{quote}

Plural \emph{i draghi}.

\begin{cousinweb}
\emph{drago} is inherited from Latin \emph{dracō} \textquotedblleft{}a serpent, a dragon\textquotedblright{} (EtymDB), from Greek \emph{drákōn}.
\end{cousinweb}

\begin{grammarlens}[title={What you've built}]
One more word for telling, explaining and describing at B1.
\end{grammarlens}

\subsection*{Your turn}
\begin{itemize}
\raggedright
  \item \emph{Say it:} \emph{il drago}
  \item \emph{Say it:} \emph{il drago}, once more
  \item \emph{Say it:} a sentence of your own with \emph{il drago}
  \item \emph{Recall:} say the Italian for the kingdom, then the Italian for to reign, then say \emph{il drago} again
\end{itemize}

\begin{tcolorbox}[breakable,colback=teal!4,colframe=teal!35!black,title={Before you move on}]
What does il drago mean? (\textquotedblleft{}The dragon\textquotedblright{}.) Say it once more.
\end{tcolorbox}
\section[la strega]{la strega — the witch}

\label{lesson:IT-C326-strega}



\begin{quote}
Before the new one: say the Italian for to reign, then the Italian for the dragon.
\end{quote}

\subsection*{You'll want to know: la strega}
\textbf{la strega} — \textquotedblleft{}the witch\textquotedblright{}.

\begin{quote}
La strega vive nel bosco.
\end{quote}

\begin{quote}
\emph{The witch lives in the woods.}
\end{quote}

Plural \emph{le streghe}.

\begin{cousinweb}
\emph{strega} is inherited from Latin \emph{strīga} \textquotedblleft{}an evil spirit\textquotedblright{} (EtymDB), built on \emph{strīx} \textquotedblleft{}a screech owl\textquotedblright{}.
\end{cousinweb}

\begin{grammarlens}[title={What you've built}]
One more word for telling, explaining and describing at B1.
\end{grammarlens}

\subsection*{Your turn}
\begin{itemize}
\raggedright
  \item \emph{Say it:} \emph{la strega}
  \item \emph{Say it:} \emph{la strega}, once more
  \item \emph{Say it:} a sentence of your own with \emph{la strega}
  \item \emph{Recall:} say the Italian for to reign, then the Italian for the dragon, then say \emph{la strega} again
\end{itemize}

\begin{tcolorbox}[breakable,colback=teal!4,colframe=teal!35!black,title={Before you move on}]
What does la strega mean? (\textquotedblleft{}The witch\textquotedblright{}.) Say it once more.
\end{tcolorbox}
\section[la fata]{la fata — the fairy}

\label{lesson:IT-C326-fata}



\begin{quote}
Before the new one: say the Italian for the dragon, then the Italian for the witch.
\end{quote}

\subsection*{You'll want to know: la fata}
\textbf{la fata} — \textquotedblleft{}the fairy\textquotedblright{}.

\begin{quote}
La fata ha aiutato la ragazza.
\end{quote}

\begin{quote}
\emph{The fairy helped the girl.}
\end{quote}

\emph{Una fata madrina} — a fairy godmother.

\begin{cousinweb}
\emph{fata} is inherited from Vulgar Latin \emph{Fāta} \textquotedblleft{}goddess of fate\textquotedblright{} (EtymDB), built on Latin \emph{fātum} \textquotedblleft{}destiny\textquotedblright{}.
\end{cousinweb}

\begin{grammarlens}[title={What you've built}]
One more word for telling, explaining and describing at B1.
\end{grammarlens}

\subsection*{Your turn}
\begin{itemize}
\raggedright
  \item \emph{Say it:} \emph{la fata}
  \item \emph{Say it:} \emph{la fata}, once more
  \item \emph{Say it:} a sentence of your own with \emph{la fata}
  \item \emph{Recall:} say the Italian for the dragon, then the Italian for the witch, then say \emph{la fata} again
\end{itemize}

\begin{tcolorbox}[breakable,colback=teal!4,colframe=teal!35!black,title={Before you move on}]
What does la fata mean? (\textquotedblleft{}The fairy\textquotedblright{}.) Say it once more.
\end{tcolorbox}
\section[il mago]{il mago — the wizard, the magician}

\label{lesson:IT-C326-mago}



\begin{quote}
Before the new one: say the Italian for the witch, then the Italian for the fairy.
\end{quote}

\subsection*{You'll want to know: il mago}
\textbf{il mago} — \textquotedblleft{}the wizard, the magician\textquotedblright{}.

\begin{quote}
Il mago ha fatto sparire il coniglio.
\end{quote}

\begin{quote}
\emph{The magician made the rabbit disappear.}
\end{quote}

Also an expert: \emph{un mago dei computer}.

\begin{cousinweb}
\emph{mago} is inherited from Latin \emph{magus} (EtymDB), from Greek \emph{mágos}, a Persian priest.
\end{cousinweb}

\begin{grammarlens}[title={What you've built}]
One more word for telling, explaining and describing at B1.
\end{grammarlens}

\subsection*{Your turn}
\begin{itemize}
\raggedright
  \item \emph{Say it:} \emph{il mago}
  \item \emph{Say it:} \emph{il mago}, once more
  \item \emph{Say it:} a sentence of your own with \emph{il mago}
  \item \emph{Recall:} say the Italian for the witch, then the Italian for the fairy, then say \emph{il mago} again
\end{itemize}

\begin{tcolorbox}[breakable,colback=teal!4,colframe=teal!35!black,title={Before you move on}]
What does il mago mean? (\textquotedblleft{}The wizard, the magician\textquotedblright{}.) Say it once more.
\end{tcolorbox}
\section[il tesoro]{il tesoro — the treasure; darling}

\label{lesson:IT-C326-tesoro}



\begin{quote}
Before the new one: say the Italian for the fairy, then the Italian for the wizard.
\end{quote}

\subsection*{You'll want to know: il tesoro}
\textbf{il tesoro} — \textquotedblleft{}the treasure; darling\textquotedblright{}.

\begin{quote}
Hanno trovato il tesoro sotto un albero.
\end{quote}

\begin{quote}
\emph{They found the treasure under a tree.}
\end{quote}

\emph{Tesoro!} — darling!

\begin{cousinweb}
\emph{tesoro} is inherited from Latin \emph{thēsaurus} \textquotedblleft{}a treasure, a store\textquotedblright{} (EtymDB), from Greek \emph{thēsaurós}.
\end{cousinweb}

\begin{grammarlens}[title={What you've built}]
One more word for telling, explaining and describing at B1.
\end{grammarlens}

\subsection*{Your turn}
\begin{itemize}
\raggedright
  \item \emph{Say it:} \emph{il tesoro}
  \item \emph{Say it:} \emph{il tesoro}, once more
  \item \emph{Say it:} a sentence of your own with \emph{il tesoro}
  \item \emph{Recall:} say the Italian for the fairy, then the Italian for the wizard, then say \emph{il tesoro} again
\end{itemize}

\begin{tcolorbox}[breakable,colback=teal!4,colframe=teal!35!black,title={Before you move on}]
What does il tesoro mean? (\textquotedblleft{}The treasure; darling\textquotedblright{}.) Say it once more.
\end{tcolorbox}
